\documentclass[11pt]{article}
\usepackage[margin=1in]{geometry}
\usepackage{booktabs}
\usepackage{enumitem}
\usepackage{xcolor}
\usepackage[hidelinks]{hyperref}
\usepackage{microtype}

\definecolor{responseblue}{RGB}{24,72,120}
\newcommand{\response}[1]{\par\noindent\textcolor{responseblue}{\textbf{Response:}} #1\par}
\newcommand{\change}[1]{\par\noindent\textcolor{responseblue}{\textbf{Revision:}} #1\par}

\title{Response to Reviewers\\
\large Manuscript HCCOM-D-26-00263\\
\textit{NeuroCursor: A Behavioral Biometric Framework for Mouse-Based User Authentication}}
\author{Shaurya Saxena, Shreehan Aalok Pathak, Hamid Abdu Mossa, and Chetan Gangireddy}
\date{}

\begin{document}
\maketitle

We thank the editor and reviewers for the careful and constructive evaluation.
We substantially revised the manuscript rather than defending claims that the
current evidence cannot support. The revision shortens the Introduction,
reorganizes Related Work around direct comparisons, states the exact input to
every model family, adds an explanatory feature-group analysis, treats the
CNN--GRU as a rejected candidate rather than a retained system choice, and
explicitly delimits what the study establishes about enrollment size and
temporal dependence. The principal changes are summarized below.

\begin{center}
\begin{tabular}{p{0.29\textwidth}p{0.62\textwidth}}
\toprule
\textbf{Issue} & \textbf{Major revision} \\
\midrule
Length and repetition & Rewrote the Introduction and Related Work; removed
duplicated architecture motivation and equations from Related Work. \\
Model inputs & Added exact neural tensor and classical vector dimensions,
including all 13 derived classical statistics. \\
Model interpretation & Added grouped coefficient analysis and a more cautious
interpretation of the ablations. \\
CNN--GRU justification & No longer recommends or ``retains'' CNN--GRU; reports
it as a tested negative result because GRU-only and classical models perform
better on key metrics. \\
Enrollment sufficiency & Clarifies that three sessions is a split-eligibility
rule, not a stable minimum; adds a nested learning curve across three splits. \\
Temporal dependence & Removes any implication that mutual information was
demonstrated; adds retrained time-permutation and low-order feature controls
for CNN--GRU and GRU-only. \\
\bottomrule
\end{tabular}
\end{center}

\section*{Reviewer 1}

\subsection*{Comment 1: Introduction and Related Work are repetitive}
The reviewer observed that both sections repeatedly motivate mouse-based
authentication and the CNN/GRU architecture, and asked that Related Work focus
on differences from prior systems.

\response{We agree. The earlier structure repeated the same motivation and
architecture rationale. We rewrote both sections. The Introduction now states
the problem, protocol, primary result, research questions, and four concise
contributions without subsections or mathematical preliminaries. Related Work
now compares task design, data, representation, strength, and limitation rather
than re-explaining the proposed model.}

\change{The revised Related Work includes Table~I, ``Positioning Relative to
Representative Mouse-Dynamics Studies,'' which directly contrasts Antal and
Fej\'er, Antal and Egyed-Zsigmond, Kuric \textit{et al.}, LT-AMouse, and
NeuroCursor. It also states that the CNN--GRU architecture is not the novelty.}

\subsection*{Comment 2: Defined features and actual model inputs are unclear}
The reviewer requested a summary of the features supplied to the neural models
and the statistics used by the classical models.

\response{We agree. The earlier temporal-channel table did not explain that
the neural and classical models use different mathematical representations of
the same source channels. The revision now states this distinction explicitly.}

\change{The new ``Exact Inputs Used by Model Family'' table reports: raw,
$d=5$, neural $128\times5$, classical 20 values; raw plus spatial and
kinematic, $d=17$, neural $128\times17$, classical 68 values; and full,
$d=27$, neural $128\times27$, classical 121 values. Classical vectors contain
mean, standard deviation, minimum, and maximum for every selected channel. The full vector
also contains 13 derived statistics: duration, path length, start-to-target
distance, path efficiency, mean/maximum speed, mean/maximum acceleration,
pause duration, click timing, acceleration energy, jerk energy, and
dimensionless jerk. The architecture section also clarifies that the CNN emits
an ordered $32\times64$ representation to the GRU.}

\subsection*{Comment 3: Additional experiments should explain model and feature behavior}
The reviewer noted that GRU-only and logistic regression outperform CNN--GRU
and that the full feature set is not consistently best, and suggested analyses
of data availability and noise-sensitive features.

\response{We agree with the interpretation. We do not present the hybrid as a
winner. We added a reproducible post-hoc grouping of all 121 standardized
logistic coefficients across the 30 refitted splits. Timing/click features have
the largest group mean absolute coefficient (0.796), followed by kinematic
(0.463) and spatial/geometric/target features (0.393). We also emphasize that
raw plus spatial and kinematic features have slightly worse AUC, FRR, and EER
than raw channels in the fixed CNN--GRU ablation. This provides evidence consistent with
noise-sensitive differentiation, while explicitly avoiding causal claims.}

\change{The Feature Ablation subsection now includes the grouped analysis and
explains why correlated standardized coefficients are descriptive rather than
causal. The Discussion identifies two plausible but unproven explanations for
the hybrid's underperformance---redundancy with engineered channels and greater
estimation difficulty in a small cohort. We additionally ran a nested logistic
learning curve across seeds 0, 1, and 2, and retrained CNN--GRU and GRU-only with
full ordered, full time-permuted, and low-order inputs at seed 42. Removing
acceleration, jerk, angular velocity, and curvature reduces hybrid EER from
7.76\% to 2.52\%, but raises GRU-only EER from 1.09\% to 7.86\%.
The low-order logistic representation also performs worse overall. These
architecture-dependent outcomes support a cautious feature-quality explanation,
not a blanket claim that derivative features are harmful. The new Results
subsection reports all conditions, including unfavorable findings.}

\section*{Reviewer 2}

\subsection*{Comment 1: The Introduction is overly long and the contributions are hard to identify}

\response{We agree. The Introduction was rewritten from four subsections into
a compact narrative. It now reports the primary repeated-split result early and
ends with four numbered contributions. Repeated explanations of passwords,
biometrics, trajectory equations, and CNN/GRU roles were removed.}

\subsection*{Comment 2: Related Work is fragmented, verbose, and unclear about recent work}

\response{We agree. Related Work is now a single comparative narrative with no
architecture equations and a five-row positioning table. The revision states
the strength and limitation of each representative study. In particular, it
acknowledges that LT-AMouse is stronger on explicit data-sufficiency analysis,
segment-length analysis, public-dataset coverage, and attack settings, whereas
NeuroCursor studies a bounded randomized target challenge and one-to-five
movement decisions.}

\subsection*{Comment 3: Does training-data scale affect accuracy, and what is the minimum enrollment data?}

\response{Yes, training scale materially affects the measured performance.
We ran new logistic learning curves across seeds 0, 1, and 2 with nested
training-session caps of 1, 3, 5, 10, 20, and all available sessions. Within
each split, validation and test sessions remain fixed, preprocessing is fitted
on the selected training sessions, and thresholds are recalibrated exclusively
on validation data. Mean strict FRR falls from 82.22\% at one session to
37.78\% at the five-session cap and 32.78\% at the full budget.}

\change{The new Results subsection includes a complete FAR/FRR/EER table and
archives actual per-user session counts. Budgets are capped by available
history: one identity has only one training session, while the others have
17--25 at the full budget. FRR is nonmonotonic across intermediate budgets and
the full-budget FRR ranges from 21.67\% to 45.00\% across the three splits.
Accordingly, the observed data do not support a stable universal minimum.
The three-session eligibility rule remains a splitting requirement. Larger
longitudinal cohorts and replicated learning curves are required to estimate
a deployment enrollment minimum.}

\subsection*{Comment 4: What justifies retaining CNN--GRU after H1 is rejected?}

\response{Nothing in the present dataset justifies choosing CNN--GRU over
GRU-only. We revised the manuscript accordingly. The hybrid is included only
as the prespecified candidate needed to report the rejected hypothesis and the
negative result. It is no longer described as the recommended or retained
system model.}

\change{The architecture section is renamed ``Candidate Neural
Architectures.'' The Results, Discussion, and Conclusion state that GRU-only
has lower FAR, FRR, and EER in the fixed neural comparison, that several
classical models are also stronger, and that the present evidence does not
support retaining the hybrid.}

\subsection*{Comment 5: Does the model capture mutual information across sampling points rather than acceleration-like features?}

\response{Architecture alone does not demonstrate cross-sample mutual
information. We therefore added a direct order-sensitivity control: retrain
both networks after independently permuting the 128 sample vectors within each
trajectory. This preserves each channel's marginal distribution and the
cross-channel relationships at each sample, while destroying chronological
order. Splits, architecture, epoch limit, early stopping, and initialization
seed are held fixed across the fresh conditions.}

\change{Fresh seed-42 GRU-only EER rises from 1.09\% with ordered full inputs
to 8.10\% after permutation, consistent with useful chronological structure.
CNN--GRU EER instead falls from 7.76\% to 4.86\%, so this control does not
support temporal-order benefit for the hybrid. Low-order feature controls are
also reported. These are new fits, distinct from the historical comparison;
stochastic fitting can yield different weights despite the same nominal seed.
One-split controls cannot establish causality or estimate mutual information.
We removed any implication that such information was demonstrated, and state
that reversal, matched-capacity controls, and replication remain necessary.}

\subsection*{Comment 6: Major revision is required}

\response{We agree and performed a substantive revision across the Abstract,
Introduction, Related Work, model-input description, architecture
interpretation, methodology scope, feature analysis, Discussion, Limitations,
and Conclusion. We also added a small reproducibility script that recomputes
the reviewer-requested input dimensions and grouped coefficient statistics
from the archived result bundle.}

\section*{Final verification}

We recompiled and visually inspected the revised manuscript and this response.
We also checked every numerical value against the archived result files,
verified that all cross-references resolve, confirmed that no new performance
result was asserted without a reproducible source, and confirmed that every
review comment is mapped to a manuscript change above.

\end{document}
